\fsection{Ej 5 Pág 125}
\fsubsection
[\textcolor{waveBlue2}{ PENSAMIENTO CRÍTICO . }]
{\textcolor{waveBlue2}{\boldit{ PENSAMIENTO CRÍTICO . }}}
{
	¿ Por qué crees que , en general , en los países más desarrollados la huella ecológica es mayor y
	el día del sobregiro se produce antes ?
}

\begin{solution}
	La relación entre el desarrollo económico y la huella ecológica se debe principalmente a los siguientes factores:

	\begin{enumerate}
		\item \textbf{Patrones de consumo:} En los países más desarrollados, la renta per cápita es mayor, lo que se traduce en un consumo intensivo de bienes, servicios y energía. Esto implica una mayor demanda de recursos naturales para producir lo que se consume (desde tecnología hasta alimentos procesados).
		\item \textbf{Dependencia de combustibles fósiles:} Estas economías suelen tener infraestructuras y sistemas de transporte (vehículos privados, aviación) que dependen fuertemente de la energía. Esto dispara la huella de carbono, que es el componente principal de la huella ecológica.
		\item \textbf{Generación de residuos y emisiones:} Los procesos industriales y el estilo de vida de "usar y tirar" generan una cantidad de residuos y emisiones de CO2 que superan la capacidad de regeneración y absorción de los ecosistemas locales.
	\end{enumerate}

	Como resultado, estos países consumen su "presupuesto ecológico" anual mucho más rápido que los países en vías de desarrollo. Por ello, su \textbf{día del sobregiro} (Earth Overshoot Day) ocurre en los primeros meses del año (febrero, marzo o abril), lo que significa que necesitarían varios planetas Tierra para mantener ese ritmo de vida de forma sostenible.
\end{solution}
